% ================================================
% 文件: 01_发射机.tex
% 章节: 第13章 发射机
% 所属部分: 发射机技术
% 版本: v1.0
% ================================================
% 本章目录结构（树状）:
% \chapter{发射机}
% ├── \section{发射机概述}
% ├── \section{发射机基本原理}
% ├── \section{高频功率放大器}
% ├── \section{调制器与解调器}
% └── \section{发射机设计与调试}
%     ├── \subsection{发射机设计要点}
%     ├── \subsection{频率合成技术}
%     ├── \subsection{发射机指标测试}
%     └── \subsection{发射机效率计算}
% ================================================

\chapter{发射机}
\label{chap:transmitter}

\section{发射机概述}
\label{sec:transmitter_intro}

发射机是一种将低频信号（如音频、视频或数字基带数据）调制到高频载波上，并通过天线以电磁波形式辐射到空间的电子设备。
发射机是无线电通信系统的核心组成部分，其电气性能与结构可靠性直接决定了通信距离、链路质量以及频谱利用效率。
在一条完整的通信链路中，发射机与接收机、传播信道共同构成闭环，其中发射机负责将待传信息"搬移"至适合空间辐射的射频频带并完成功率放大。

从系统功能看，发射机的主要任务可概括为三点：其一是产生频率准确、稳定且相位噪声足够低的载波；其二是按预定方式将基带信息加载到载波上（即调制）；其三是将已调信号放大到额定射频功率，使之经天线辐射后能在接收端获得足够的解调信噪比。
现代发射机普遍采用频率合成、线性功放与数字预失真等技术，以兼顾频谱纯度、效率与多制式兼容。

发射机可按多种维度分类。
按输出功率等级可分为小功率（毫瓦至数瓦，常用于遥控、数传与便携设备）、中功率（数瓦至数百瓦，常用于基地台与广播电视差转）与大功率（数百瓦至数十千瓦，用于广播与干线通信）发射机；
按工作频段可分为长波、中波、短波、超短波与微波发射机，不同频段对应不同的传播方式与天线形式；
按调制方式可分为调幅（AM）、调频（FM）、单边带（SSB）、调相（PM）及各类数字调制发射机。

\begin{itemize}
    \item \textbf{按功率等级}：小功率发射机、中功率发射机、大功率发射机
    \item \textbf{按工作频段}：长波、中波、短波、超短波、微波发射机
    \item \textbf{按调制方式}：AM发射机、FM发射机、SSB发射机、数字调制发射机
\end{itemize}

\begin{table}[htbp]
\centering
\caption{发射机按调制方式的分类对照}
\label{tab:tx_classification}
\begin{tabular}{|p{3.2cm}|p{4.3cm}|p{5cm}|}
\hline
\textbf{调制方式} & \textbf{典型应用} & \textbf{主要特点} \\
\hline
调幅（AM） & 中短波广播 & 电路简单，带宽窄，抗衰落差 \\
\hline
调频（FM） & 超短波广播、模拟对讲 & 抗噪声强，占用带宽较宽 \\
\hline
单边带（SSB） & 短波通信、船岸通信 & 功率利用率高，带宽约一半 \\
\hline
数字调制 & 移动通信、数据链路 & 频谱效率高，需线性功放 \\
\hline
\end{tabular}
\end{table}

发射机品质的优劣通常用输出功率、频率稳定度、调制准确度、邻道泄漏比（ACLR）、占用带宽与杂散辐射等指标衡量，这些指标既是工程设计的目标，也是型号核准与电磁兼容（EMC）检测的依据。
有效辐射功率（ERP）将发射功率与天线增益联系起来：
\begin{equation}
P_{\mathrm{ERP}} = P_{\mathrm{tx}} \times G_{\mathrm{tx}},
\end{equation}
其中 $P_{\mathrm{tx}}$ 为天线端口发射功率，$G_{\mathrm{tx}}$ 为天线相对于半波振子的增益（倍数）。

现代发射机正加速向软件定义与绿色化演进：以现场可编程门阵列为调制/波形生成核心，配合高速数模变换与数字预失真（DPD），可在同一射频前端支持多种制式与带宽；同时借助 Doherty、包络跟踪等高效功放技术降低功耗与散热需求。
各国无线电管理机构对发射模板、占用带宽、邻道泄漏与杂散均设有明确限值，发射机设计须以"合规模板"为硬约束，而非仅追求峰值功率。

\section{发射机基本原理}
\label{sec:transmitter_principle}

发射机的基本工作过程可概括为：由信号源产生原始低频（基带）信号，经调制器将其加载到高频载波上，再经若干级变频与放大、滤波后送至天线辐射。
对采用超外差或正交上变频架构的现代发射机而言，典型链路为：基带信源 $\rightarrow$ 调制/数模变换 $\rightarrow$ 中频（IF）处理 $\rightarrow$ 上变频至射频（RF）$\rightarrow$ 射频驱动放大 $\rightarrow$ 末级功率放大 $\rightarrow$ 谐波/带通滤波 $\rightarrow$ 环形器/天线开关 $\rightarrow$ 天线。
各功能级的增益、线性度与相位噪声沿链路逐级叠加，因此载波纯度主要由前端频率合成器决定，而带外杂散与邻道泄漏则更多受末级功放非线性影响。

信号源提供待传输的信息，可能是模拟音频/视频，也可能是经信源编码与信道编码的数字比特流。
调制器按选定制式将基带信息映射为已调波形；对于调幅，已调包络随调制信号变化，调幅指数定义为
\begin{equation}
m = \frac{A_{\max}-A_{\min}}{A_{\max}+A_{\min}} = \frac{k_a U_{\Omega}}{U_c},
\end{equation}
其中 $U_c$ 为载波幅度，$U_{\Omega}$ 为调制信号幅度，须满足 $m\le 1$ 以避免过调幅失真。
对于调频，瞬时频偏 $\Delta f$ 与调制幅度成正比，按卡森（Carson）公式其近似占用带宽为
\begin{equation}
B_T \approx 2(\Delta f + f_m),
\end{equation}
式中 $f_m$ 为最高调制频率。

上变频将中频已调信号搬移到射频工作频率，常用方案为
\begin{equation}
f_{\mathrm{RF}} = f_{\mathrm{LO}} \pm f_{\mathrm{IF}},
\end{equation}
即通过混频器将本振（LO）与中频（IF）相加或相减得到射频。
该过程会同时产生和频与差频两个分量，必须用其后滤波器抑除不想要的镜像与杂散。
经过变频的信号再经驱动级与末级功率放大器提升幅度，最后由带通/低通滤波网络滤除谐波后馈至天线。

发射机的核心指标包括输出功率、总效率、频率准确度与稳定度、相位噪声、调制精度（调幅指数/调频频偏/误差矢量幅度 EVM）、邻道功率泄漏比（ACLR）以及占用带宽与杂散辐射。
它们彼此间往往存在折衷，例如提高功放效率常以牺牲线性度为代价，工程上需针对业务场景统筹取舍。

除指标折衷外，发射机还需建立自我保护闭环：ALC（自动电平控制）环路实时采样输出幅度并反馈调节驱动增益，抑制电源与温度引起的功率漂移；驻波（VSWR）检测在天线失配时降低驱动或切断功放，避免反射功率损坏末级；过温与过流保护则在散热异常时降额或关机。
良好的保护策略能在不影响正常指标的前提下显著提升现场可靠性，是工程发射机与实验电路的显著区别之一。

\section{高频功率放大器}
\label{sec:rf_power_amplifier}

高频功率放大器是发射机中将已调信号放大到额定射频功率的关键部件，其输出经滤波后直接驱动天线。
与接收通道的小信号放大器不同，功放工作于大信号非线性区，设计目标是在给定输出功率、线性度与带宽约束下尽可能提高能量转换效率并降低热耗。
功放的效率与晶体管导通角密切相关，按导通角（一个射频周期内集电极/漏极电流导通的电角度 $\theta$）可分为 A、B、AB、C 及开关类（D、E、F）等若干类别。

A 类（甲类）放大器在整个周期内均导通（$\theta = 360^\circ$），线性最好，但理论最高集电极效率仅约 $50\%$（实际约 $25\%\sim 50\%$）；
B 类（乙类）仅在半周导通（$\theta = 180^\circ$），理论效率约 $78.5\%$，但存在交越失真，常用于推挽；
AB 类是二者的折衷；
C 类（丙类）导通角小于 $180^\circ$，效率可达 $80\%\sim 90\%$，但输出含丰富谐波，必须经调谐回路选频，仅适用于恒包络（如 FM、CW）发射；
D/E 类以开关方式工作，理论效率可超过 $90\%$，多用于音频与较低射频功率合成。

常用效率定义如下。集电极（漏极）效率：
\begin{equation}
\eta_c = \frac{P_{\mathrm{RF}}}{P_{\mathrm{dc}}} = \frac{P_{\mathrm{RF}}}{P_{\mathrm{RF}}+P_c},
\end{equation}
其中 $P_{\mathrm{RF}}$ 为射频输出功率，$P_{\mathrm{dc}}$ 为直流输入功率，$P_c$ 为晶体管耗散。
更能反映"可用线性功率"的功率附加效率（PAE）为
\begin{equation}
\mathrm{PAE} = \frac{P_{\mathrm{RF}}-P_{\mathrm{in}}}{P_{\mathrm{dc}}}\times 100\%,
\end{equation}
$P_{\mathrm{in}}$ 为射频输入驱动功率。显然，当驱动功率远小于输出功率时，PAE 趋近于集电极效率。

实际功放设计还需关注：偏置电路的稳定性与温度补偿、输入/输出匹配网络（常用 L 型、$\pi$ 型、T 型或传输线变压器）以实现共轭匹配并兼顾谐波抑制、热设计与散热基板、以及负载牵引（Load-Pull）下的安全工作区。
匹配网络的品质因数 $Q$ 应适度，过高会压缩带宽、过低则滤波不足。
此外，末级功放通常采用 ALC（自动电平控制）稳定输出幅度，并借助环形器或隔离器降低负载失配对功放的牵引。

对于线性数字发射机，功放常工作于回退区以保线性，效率因而下降，需引入线性化技术补偿：数字预失真（DPD）在基带对信号预加逆失真、抵消功放非线性；前馈与笛卡尔反馈用于高线性场景；包络跟踪（ET）则让供电随包络变化，使功放在高峰均比信号下仍保持高效。
线性化与高效功放的协同，是现代宽带发射机设计的核心难点，直接决定 ACLR 与误差矢量幅度（EVM）是否满足模板。

\begin{table}[htbp]
\centering
\caption{高频功率放大器工作类别对照}
\label{tab:pa_classes}
\begin{tabular}{|p{2.2cm}|p{3cm}|p{3.6cm}|p{4cm}|}
\hline
\textbf{类别} & \textbf{导通角} & \textbf{理论/典型效率} & \textbf{适用场景} \\
\hline
A 类 & $360^\circ$ & 约 $25\%\sim 50\%$ & 高线性小功率 \\
\hline
B 类 & $180^\circ$ & 约 $78.5\%$ & 推挽射频放大 \\
\hline
C 类 & $<180^\circ$ & 约 $80\%\sim 90\%$ & 恒包络 FM/CW \\
\hline
D/E 类 & 开关 & $>90\%$ & 高效开关功放 \\
\hline
\end{tabular}
\end{table}

\section{调制器与解调器}
\label{sec:modulator_demodulator}

调制器用于将低频信号加载到高频载波上，使信息得以在射频信道中传输；解调器则完成逆过程，在接收端从已调波中恢复原始信息。
尽管本节立足于发射机，但调制与解调在原理上成对出现，其性能指标（如调制线性度、解调信噪比）应统一考量。
常见调制方式可按载波受控参数分为调幅、调频、调相三大类，以及在此基础上发展的各类数字调制。

调幅（AM）使载波幅度随调制信号变化，单音调制时已调波可写为
\begin{equation}
s_{\mathrm{AM}}(t)=A_c\bigl[1+m\cos(\omega_m t)\bigr]\cos(\omega_c t),
\end{equation}
其频谱由载频与上下边带组成，带宽为基带最高频率的两倍。
调频（FM）使瞬时频率随调制信号偏移，调相（PM）则使瞬时相位偏移，二者统称角度调制，具有恒包络、抗幅度噪声强的优点，但占用带宽较宽。
数字调制以离散符号改变载波的幅度、频率或相位，如幅移键控（ASK）、频移键控（FSK）、相移键控（PSK）及正交幅度调制（QAM）；其中 QAM 同时改变幅度与相位，频谱效率高，但对功放线性度要求严格，常需配合数字预失真（DPD）与线性化技术。

调制器的设计需综合考虑调制指数（或频偏、星座映射）、已调信号带宽、线性度与载波抑制比。
例如 SSB 发射需通过滤波法或相移法抑除一个边带以节省功率与带宽；而线性数字发射机则要求功放在回退（Back-off）工作点保持足够线性，用误差矢量幅度（EVM）衡量解调质量，典型要求为 EVM 优于数个百分点。
解调方面，相干解调需本地恢复同频同相载波，非相干解调（如包络检波、限幅鉴频）则结构更简单但性能略低。
调制方式的选择本质上是带宽效率与功率效率的折衷。

在线性发射机中，调制器与功放需协同设计：恒包络调制（如 FM、GMSK）允许功放工作于高效率非线性区；而振幅敏感的 QAM/SSB 则要求功放在邻近饱和区仍保持线性，故须配合预失真与适当回退。
现代收发信机多采用正交（I/Q）调制架构，由两路正交载波合成任意波形，其性能受 I/Q 失调、本振泄漏与相位噪声制约，需通过校准与滤波保证调制纯度。

\section{发射机设计与调试}
\label{sec:transmitter_design}

\subsection{发射机设计要点}

发射机设计是一项在性能、效率、成本与可靠性之间反复折衷的系统工程，需在方案阶段即明确工作频段、输出功率、调制制式与频谱模板等约束。主要设计要点如下：

\begin{itemize}
    \item \textbf{频率合成}：提供频率准确、相位噪声低的载波源，可采用晶体振荡器、锁相环（PLL）或直接数字频率合成器（DDS），并按频谱模板要求分配各链路级的相位噪声预算。
    \item \textbf{功率放大}：选用合适类别的功率放大器，在输出功率、线性度与效率间折衷；合理设计散热与热沉，确保满功率长时间工作不过温。
    \item \textbf{调制电路}：保证调制精度与线性度，数字调制需考虑预失真与 IQ 失调校正。
    \item \textbf{滤波网络}：在功放输出端设置低通/带通滤波以抑制谐波与杂散，满足电磁兼容（EMC）与型号核准要求。
    \item \textbf{电源供应}：采用稳定、高效的电源并抑制纹波与噪声，射频地与功率地合理分区以降低互扰。
    \item \textbf{保护电路}：设置过压、过流、过热及驻波（VSWR）保护，防止失配或故障损坏末级功放。
\end{itemize}

除上述要点外，还应进行链路预算与电磁兼容（EMC）设计：前者用于核算从功放到天线、经传播至接收端的可用信噪比，后者通过屏蔽、滤波与接地降低对内对外的干扰。
对于可搬移设备，体积、质量与功耗亦是关键约束。

\subsection{频率合成技术}

频率合成器用于产生稳定、准确且可捷变切换的载波频率，是发射机频率基准的核心。主流技术包括晶体振荡器、锁相环（PLL）与直接数字频率合成器（DDS），现代设备常将三者组合使用。

\begin{itemize}
    \item \textbf{晶体振荡器}：以石英晶体的压电谐振提供高稳定度、低相噪的固定频率参考，常作为 PLL/DDS 的基准源；其频率准确度受温度影响，温补晶振（TCXO）与恒温晶振（OCXO）可进一步改善。
    \item \textbf{锁相环（PLL）}：通过反馈环路将压控振荡器（VCO）锁定到参考频率的整数或分数倍，兼具宽频覆盖与较好相噪，是实现频率捷变的主要手段。
    \item \textbf{直接数字频率合成器（DDS）}：以相位累加器在数字域生成波形再经数模变换输出，频率分辨率极高、切换速度快、相位连续，但输出带宽与杂散抑制受限，常作基准或中频源。
\end{itemize}

PLL 的基本组成包括参考振荡器、鉴相器（PD）、环路滤波器（LF）与压控振荡器（VCO），其输出频率为
\begin{equation}
f_{\mathrm{out}} = f_{\mathrm{ref}}\left(N+\frac{B}{F}\right) \quad\text{或}\quad f_{\mathrm{out}} = f_{\mathrm{ref}}\times\frac{N}{R},
\end{equation}
其中 $f_{\mathrm{ref}}$ 为参考频率，$N$ 为分频比，$B/F$ 为小数分频修正项（或 $R$ 为参考分频比）。
环路带宽需在"快速锁定"与"低相噪/低杂散"之间折衷：较窄环路可有效抑制 VCO 近端噪声，较宽环路则锁定更快。
DDS 输出频率由采样时钟 $f_s$ 与相位增量 $\Delta\varphi$ 决定：
\begin{equation}
f_{\mathrm{DDS}} = \frac{\Delta\varphi}{2^{M}}\,f_s,
\end{equation}
$M$ 为相位累加器位数，故频率分辨率可达 $f_s/2^{M}$ 量级，适合需要微步长扫频或快速跳频的场合。

\subsection{发射机指标测试}

发射机研制与生产需进行系统调试与指标测试，以验证其符合设计规范与无线电管理规定。典型测试项目包括：

\begin{itemize}
    \item \textbf{功率测量}：用功率计或频谱仪测输出功率与总效率，注意源/负载驻波引起的测量误差。
    \item \textbf{频率校准}：验证载波频率准确度与短期/长期稳定度，常需频率计或频谱仪频偏测量。
    \item \textbf{调制参数}：测量调幅指数、调频频偏或数字调制的 EVM、星座图与码域误差。
    \item \textbf{谐波抑制}：测二次、三次谐波相对载波的电平（谐波抑制比），通常要求优于 $-30\sim -40\ \mathrm{dBc}$。
    \item \textbf{杂散辐射}：检测非谐波关系的杂散分量，须低于法规限值。
    \item \textbf{邻道功率}：测邻道功率泄漏比（ACLR），评估对相邻频道的干扰。
    \item \textbf{占用带宽}：按所需带宽（OBW）或 $99\%$ 能量带宽判定信号频谱占用。
\end{itemize}

表\ref{tab:tx_test} 汇总了主要测试项目、常用仪器与关注要点。测试中应使用经过校准的衰减器与定向耦合器，避免大功率反向损坏仪表，并注意频谱模板（Spectrum Mask）合规判定。

\begin{table}[htbp]
\centering
\caption{发射机主要测试项目}
\label{tab:tx_test}
\begin{tabular}{|p{3cm}|p{4cm}|p{5cm}|}
\hline
\textbf{测试项目} & \textbf{常用仪器} & \textbf{关注要点} \\
\hline
输出功率 & 功率计、频谱仪 & 额定功率、总效率 \\
\hline
频率准确度 & 频率计、频谱仪 & 频偏、稳定度 \\
\hline
调制参数 & 调制分析仪 & 调幅指数/频偏/EVM \\
\hline
谐波/杂散 & 频谱分析仪 & 谐波抑制比、杂散电平 \\
\hline
邻道功率 & 频谱分析仪 & ACLR、频谱模板 \\
\hline
\end{tabular}
\end{table}

\subsection{发射机效率计算}

发射机的总效率反映直流电能转化为射频输出能量的能力，是散热设计与电源容量规划的基础。总效率定义为
\begin{equation}
\eta_{\mathrm{total}} = \frac{P_{\mathrm{out}}}{P_{\mathrm{dc}}},
\end{equation}
其中 $P_{\mathrm{out}}$ 为射频输出功率，$P_{\mathrm{dc}}$ 为整机的直流输入功率（含功放、驱动、频率合成、调制与电源自身损耗）。

末级功率放大器效率可单独表示为
\begin{equation}
\eta_{\mathrm{PA}} = \frac{P_{\mathrm{RF}}}{P_{\mathrm{dc-PA}}},
\end{equation}
其中 $P_{\mathrm{dc-PA}}$ 为功放级直流功耗。
整机效率同时受电源变换效率 $\eta_{\mathrm{PS}}$、调制/中频链路效率及功放效率共同制约，工程上常将总效率近似写为各环节效率之积：
\begin{equation}
\eta_{\mathrm{total}} \approx \eta_{\mathrm{PA}}\times \eta_{\mathrm{PS}} \times \eta_{\mathrm{other}}.
\end{equation}

例如，当功放输出 $P_{\mathrm{RF}}=50\ \mathrm{W}$、直流输入 $P_{\mathrm{dc-PA}}=110\ \mathrm{W}$ 时，$\eta_{\mathrm{PA}}\approx 45.5\%$；若电源效率约 $85\%$、其余链路损耗约 $90\%$，则整机能效约 $35\%$。
由此可见，提高功放效率（如采用 Doherty、包络跟踪或开关类功放）与高效电源是降低热耗、缩小体积的关键途径。
散热设计需依据最大连续功耗 $P_c=P_{\mathrm{dc}}-P_{\mathrm{out}}$ 选择热沉与风冷/液冷方案，保证结温在安全范围内。
